\chapter{Penjumlahan dengan Menyimpan}\label{ch:g2:addition}

Ketika satuan sebuah \omterm{def:g1:addition:def}{jumlah} melewati sembilan, sepuluh di antaranya
terikat menjadi satu puluhan baru --- dan ikatan kecil itu, yang
dipindahkan ke kolom berikutnya, adalah \emph{simpanan} yang
terkenal itu. Bab ini menjelaskan simpanan dengan lidi, lalu melatih
cara bersusun.

\section{Mengapa kita menyimpan}

\begin{example}[Mengikat sepuluh satuan]\label{ex:g2:addition:bundle}
$27 + 15$: gabungkan satuannya, $7 + 5 = 12$ satuan. Tetapi sepuluh
satuan menjadi satu puluhan! Ikatlah sepuluh di antaranya: tersisa
$2$ satuan lepas dan satu puluhan baru untuk disimpan. Sekarang
puluhannya: $2 + 1 + 1 = 4$ puluhan. Jadi $27 + 15 = 42$.
\end{example}

\begin{omfigure}
\begin{tikzpicture}[scale=0.75]
  % 12 satuan: sepuluh di antaranya diikat menjadi satu ikat, dua tersisa lepas
  \foreach \i in {0,...,9} \draw[omThm, very thick]
    (0.24*\i,0) -- (0.24*\i,1.7);
  \draw[omExo, thick, rounded corners=3pt]
    (-0.24,0.62) rectangle (2.4,1.05);
  \node[below, font=\small] at (1.08,-0.25) {$10$ diikat};
  \foreach \i in {0,1} \draw[omDef, very thick]
    (4.0+0.4*\i,0) -- (4.0+0.4*\i,1.7);
  \node[below, font=\small] at (4.2,-0.25) {$2$ tersisa};
  \node[right, font=\small, align=left] at (5.5,0.85)
    {$7 + 5 = 12$: satu puluhan dan $2$ satuan.\\
     Ikatan itulah simpanannya.};
\end{tikzpicture}

{\small Simpanan adalah sebuah ikatan: sepuluh satuan lepas menjadi
satu puluhan, yang dipindahkan ke kolom puluhan.}
\end{omfigure}

\section{Cara bersusun}

\begin{method}[Penjumlahan bersusun]\label{met:g2:addition:column}
\begin{enumerate}
  \item Tulislah bilangannya satu di bawah yang lain, satuan di bawah
        satuan, puluhan di bawah puluhan;
  \item jumlahkan satuannya; jika hasilnya melewati $9$, tulislah
        angka satuannya lalu simpan angka $1$ kecil di atas puluhan;
  \item jumlahkan puluhannya (beserta simpanan), lalu ratusannya;
  \item periksalah dengan taksiran: bulatkan setiap bilangan lalu
        jumlahkan kira-kira.
\end{enumerate}
\end{method}

\begin{example}\label{ex:g2:addition:worked}
Hitunglah $268 + 156$:
\[
\begin{array}{r}
{\scriptstyle 1}\ {\scriptstyle 1}\ \phantom{0} \\[-3pt]
2\,6\,8 \\
+\ 1\,5\,6 \\
\hline
4\,2\,4
\end{array}
\]
Satuan: $8 + 6 = 14$: tulis $4$, simpan $1$. Puluhan:
$6 + 5 + 1 = 12$: tulis $2$, simpan $1$. Ratusan:
$2 + 1 + 1 = 4$. Periksa dengan taksiran: kira-kira
$270 + 160 = 430$, dekat dengan $424$. \checkmark
\end{example}

\section{Menjumlahkan di kepala}

\begin{example}[Membuat puluhan penuh]\label{ex:g2:addition:mental}
Untuk \omterm{def:g1:addition:def}{jumlah} yang kecil, lompatlah dulu ke puluhan penuh:
\[
48 + 6 = 48 + 2 + 4 = 50 + 4 = 54 .
\]
Jumlahkan puluhan dan satuan secara terpisah:
$35 + 23 = (30 + 20) + (5 + 3) = 50 + 8 = 58$. Dan ada trik untuk
bilangan yang hampir seratus: $99 + 46 = 100 + 46 - 1 = 145$.
\end{example}

\begin{example}[Sebuah soal kecil]\label{ex:g2:addition:problem}
``Perpustakaan punya $185$ buku; $57$ buku baru datang. Berapa buku
sekarang?'' Menjumlahkan: $185 + 57 = 242$ (satuan $5 + 7 = 12$,
simpan\,\dots). Perpustakaan itu punya $242$ buku.
\end{example}

\section{Latihan}

\begin{exercise}[$\star$]\label{exo:g2:addition:1}
Hitunglah (tanpa simpanan): $34 + 25$; \quad $126 + 43$; \quad
$507 + 82$.
\end{exercise}

\begin{exercise}[$\star$]\label{exo:g2:addition:2}
Hitunglah dengan simpanan: $47 + 8$; \quad $56 + 27$; \quad $65 + 19$.
\end{exercise}

\begin{exercise}[$\star$]\label{exo:g2:addition:3}
Hitunglah secara bersusun: $268 + 145$; \quad $374 + 158$; \quad
$429 + 293$.
\end{exercise}

\begin{exercise}[$\star$]\label{exo:g2:addition:4}
Hitunglah di kepala dengan membuat puluhan penuh: $38 + 5$; \quad
$67 + 6$; \quad $54 + 8$.
\end{exercise}

\begin{exercise}[$\star$]\label{exo:g2:addition:5}
Hitunglah dengan menjumlahkan puluhan dan satuan secara terpisah:
$42 + 36$; \quad $53 + 24$; \quad $61 + 28$.
\end{exercise}

\begin{exercise}[$\star$]\label{exo:g2:addition:6}
Taksirlah dulu (bulatkan ke puluhan atau ratusan terdekat), lalu
hitunglah dengan tepat: $58 + 34$; \quad $296 + 187$.
\end{exercise}

\begin{exercise}[$\star$]\label{exo:g2:addition:7}
``Ani punya $128$ kelereng, Tono punya $95$. Berapa banyaknya kalau
digabung?'' Tulislah penjumlahannya dan kalimat jawabannya.
\end{exercise}

\begin{exercise}[$\star$]\label{exo:g2:addition:8}
Lengkapilah satuan yang kosong: $3\,?\, + 27 = 64$ (carilah angka
satuan yang hilang dari bilangan pertama).
\end{exercise}

\begin{exercise}[$\star$]\label{exo:g2:addition:9}
Dua kelas pergi menonton pertunjukan: $27$ dan $28$ anak, ditambah
$6$ orang dewasa. Berapa orang seluruhnya? (Jumlahkan kedua kelas dulu.)
\end{exercise}

\begin{exercise}[$\star\star$]\label{exo:g2:addition:10}
Carilah tiga bilangan pada daftar $27$, $33$, $40$, $65$ yang
berjumlah $100$. (Carilah dulu dua bilangan yang satuannya sahabat sepuluh.)

\end{exercise}
